\begin{Problem}[NAND and XOR]
	Let us represent Boolean values by $\text{FALSE}=0$ and $\text{TRUE}=1$. As activation function let us choose the step function
	\[
	\sigma(x)=
	\begin{cases}
		1 & \text{if } x\ge 0,\\
		0 & \text{if } x<0.
	\end{cases}
	\]
	
	\begin{enumerate}
		\item[(a)] Find suitable weights $\vec w=(w_1,w_2)$ and a bias $b$ so that a single perceptron emulates a NAND gate.
		
		\item[(b)] Show that it is impossible to realize a XOR gate with a single perceptron.
	\end{enumerate}
\end{Problem}
\begin{proof}
	\begin{parts}
		\item A NAND gate outputs 1 in all cases unless both inputs are 1, in which case it returns 0. A perceptron outputs
		\[f(x_1,x_2) = \sigma(w_1 x_1+w_2x_2 + b).\]
		Substituting in 1 and 0 for $x_1$ and $x_2$ accordingly, and using the properties of the step function, we see that we are demanding that
		\begin{align*}
			b &\ge  0\\
			w_1+w_2 + b&< 0\\
			w_1 + b &\ge 0\\
			w_2 + b &\ge 0 
		\end{align*}
	From comparing the second inequality with the third and fourth, we see that $w_1<0$ and $w_2 < 0$. The first inequality obviously says $b\ge 0$. However, we can see that there is no unique solution; for example, we have at least a constant prefactor degree of freedom on all of the parameters.
	
	Fixing the signs of all the parameters as defined above, the only constraints are
	\begin{align*}
		|w_1|, |w_2|&\le |b|\\
		|w_1|+|w_2|&> |b|
	\end{align*}
An example solution is $b=3$, $w_1=w_2=-2$.
\item An XOR gate outputs 1 if either input is TRUE, and 0 otherwise (including when both are TRUE). Substituting, we find that
\begin{align*}
	w_1 + b &\ge 0\\
	w_2 + b &\ge 0\\
	w_1+w_2+b &<0\\
			b &< 0
\end{align*}
Now, the last inequality clearly states $b<0$. Comparing the first, second and third equations again, we find that $w_1<0$ and $w_2<0$. However, combining this with the fourth equation, it is now impossible for $w_1+b$ to be positive.\qedhere
	\end{parts}
\end{proof}
\begin{Problem}[Gradient Method --- Steepest Descent]
	Since learning in artificial neural networks can be viewed as some kind of optimization problem, this exercise aims to make you familiar with the method of steepest descent.
	
	Let us first consider a parabola $f(x)=\frac12 ax^2+bx+c$. Choosing an arbitrary starting value $x_0\in\mathbb R$, let us compute the iteration $x_0 \to x_1 \to x_2 \to \dots$ through updates of the form
	\[
	x_k \to x_{k+1}:=x_k-\alpha f'(x_k),
	\]
	where $\alpha\in\mathbb R$ is a constant.
	
	\begin{enumerate}
		\item[(a)] Find the range for $\alpha$ where $\{x_k\}$ converges to the minimum of the function.
		
		\item[(b)] Determine the optimal value for $\alpha$ so that the iteration reaches the minimum already after a single iteration.
		
		\item[(c)] Use this method to find the minimum of
		\[
		f(x)=\cosh(x),
		\]
		starting with $x_0=1$. Choose the optimal $\alpha$ and carry out the computation numerically.
		
		\item[(d)] Now let us be more general and consider a function defined on a vector space. To this end, let $M$ be a real symmetric $n\times n$ matrix with nonzero eigenvalues and let the function $f:\mathbb R^n\to\mathbb R$ be given by
		\[
		f(\vec x)=\frac12(\vec x-\vec a)^T M^2 (\vec x-\vec a).
		\]
		
		Rewrite $f(\vec x)$ in the form
		\[
		f(\vec x)=\frac12 \vec x^T A \vec x+\vec b\cdot \vec x+c
		\]
		and determine $A,\vec b,c$.
		
		\item[(e)] Prove that $f(\vec x)$ has a unique global minimum at $\vec x=\vec a$.
		
		\item[(f)] The iteration rule has now to be generalized as follows:
		\[
		\vec x_k \to \vec x_{k+1}:=\vec x_k-\alpha \nabla f(\vec x_k).
		\]
		
		Compute the gradient $\nabla f(\vec x_k)$ and determine the optimal step size.
	\end{enumerate}
\end{Problem}
\begin{proof}
	We compute the minimum point:
	\[f'(x) = ax + b\]
	and thus the minimum point $x_\text{min}$ is given by
	\[x_\text{min} = -\frac ba\]
	We also note that the explicit form of the iteration is given by
	\[x_{k+1} := x_k - \alpha (ax_k + b) = (1-\alpha a)x_k + \alpha b\]
	\begin{parts}
		\item We compute the distance from the minimum. Defining $\Delta_k = x_k - x_\text{min}$, we find the following recursion relation
		\begin{align*}
			\Delta_{k+1} &= x_{k+1} - x_\text{min}\\
			&= \Delta_k - \alpha (a x_k + b)\\
			&= \Delta_k - \alpha[a(\Delta_k+x_\text{min}) + b]\\
			&= \Delta_k - \alpha a \Delta_k\\
			&= \Delta_k(1-\alpha a)
		\end{align*}
	from which we can easily prove by induction that
	\[\Delta_k = (1-\alpha a)^k \Delta_0\]
	Now, we see that the distance from the minimum either increases or decreases monotonically, a consequence of the convexity of the function. For the $\Delta_k$s to decrease, we need $|1-\alpha a|<1$, and thus
	\[0<\alpha < \frac 2a\]
	\item From the previous expression, we require that $1-\alpha a = 0$, and thus
	\[\alpha = \frac 1a\]
	\item We note that $f'(x)=\sinh x$, and we know that the minimum point is $x_\text{min}=0$. For the optimal $\alpha$, we would require that $1 - \alpha\sinh 1 = 0$, and thus $\alpha = 1/\sinh 1 = \frac{2e}{e^2-1}$. However, we can also compute a more meaningful locally optimal $\alpha$ by expanding $\cosh$ locally. Since
	\[\cosh x\approx 1 + \frac{x^2}{2} + O(x^4),\]
	by comparison with the result from part (b), we can read off $a=1$ and thus a locally optimal $\alpha$ is given by $\alpha = 1$.
	
	We perform the iteration in Python
	\begin{lstlisting}[language=Python]
		x = 1
		max_iter = 20
		alpha = 0.1
		xlst = [x]
		for i in range(max_iter):
			x -= alpha*np.sinh(x)
			xlst.append(x)\end{lstlisting}
		Performing this for $\alpha=0.1$, 0.5 and the optimal $\alpha$ determined earlier, we find the following convergence plot
		
		\begin{center}
		\begin{tikzpicture}
			\begin{axis}[
				width=12cm,
				height=8cm,
				xlabel={$x$},
				ylabel={$y$},
				grid=major,
				legend pos=north east
				]
				
				% First line
				\addplot[
				color=blue,
				mark=*,
				thick
				]
				table[
				col sep=comma,
				x=xvals,
				y=0.2
				]{blatt7_data.csv};
				
				\addlegendentry{$\alpha=0.2$}
				
				% Second line
				\addplot[
				color=red,
				mark=square*,
				thick
				]
				table[
				col sep=comma,
				x=xvals,
				y=1
				]{blatt7_data.csv};
				
				\addlegendentry{$\alpha=1$}
				
								% Second line
				\addplot[
				color=black,
				mark=triangle*,
				thick
				]
				table[
				col sep=comma,
				x=xvals,
				y=optimal
				]{blatt7_data.csv};
				\addlegendentry{$\alpha=\frac{2e}{e^2-1}$}
			\end{axis}
		\end{tikzpicture}
	\end{center}
	\item We have
	\begin{align*}
		f(\vec{x}) &= \frac 12 (\vec x - \vec a)^T \mathbf{M}^2(\vec x - \vec a)\\
		&= \frac 12 (\vec x - \vec a)^T(\mathbf{M}^2 \vec x - \mathbf{M}^2\vec a)\\
		&= \frac 12 \left[(\vec x)^T \mathbf{M}^2 \vec x - (\vec x)^T \mathbf{M}^2\vec a - (\vec a)^T \mathbf{M}^2 \vec x + (\vec a)^T \mathbf{M}^2\vec a\right] \\
		&= \frac 12 \left[(\vec x)^T \mathbf{M}^2 \vec x - (\vec a)^T \mathbf{M}^2\vec x - (\vec a)^T \mathbf{M}^2 \vec x + (\vec a)^T \mathbf{M}^2\vec a\right] \\
		&= \frac 12 (\vec x)^T \mathbf{M}^2 \vec x - (\mathbf{M}^2 \vec a)\cdot \vec x + \frac 12 (\vec a)^T \mathbf{M}^2\vec a
	\end{align*}
from which we can identify the various parameters
\begin{align*}
	\mathbf{A} &= \mathbf{M}^2\\
	\vec b &= \mathbf{M}^2\vec{a}\\
	c &= \frac 12 \vec a^T \mathbf{M}^2\vec a
\end{align*}
\item We use the original representation of $f(\vec{x})$. The key idea here is that $\mathbf{M}^2$ is positive and symmetric. We can thus diagonalise $\mathbf{M}^2$ as
\[\mathbf{M}^2 = \sum_i \epsilon_i^2 \vec v_i \vec v_i^T\]
where the $\epsilon_i$ are the eigenvalues of $\mathbf{M}$. Again, note that since $\mathbf{M}$ is real symmetric, the $\epsilon_i^2$ are all positive. Substituting, we find that
\begin{align*}
	f(\vec x) &= \frac 12 \sum_i \epsilon_i^2(\vec x - \vec a)^T \vec v_i \vec v_i^T(\vec x - \vec a) \\
	&= \frac 12 \sum_i \epsilon_i^2 |\vec v_i^T(\vec x - \vec a)|^2\\
	&\ge 0
\end{align*}
Thus, the range of $f$ is bounded from below. When $\vec x - \vec a = 0$, it is obvious that $f=0$. Thus $f$ has a minimum. Now we show that this minimum is unique. Suppose $\vec x - \vec a \neq 0$. Since the $\vec v_i$s form an orthogonal basis by the spectral theorem, it follows that $\vec v_i^T (\vec x - \vec a)\neq 0$ for at least one $i$, forcing $f$ to be positive. Thus $\vec x=\vec a$ is the only global minimum.
\item We must first differentiate the function. Summing over repeated indices (so $\mathbf{M}_{ij}x_j$ has an implicit sum over $j$), we have
	\begin{align*}
		\partial_k f(\vec x) &= \frac{1}{2}\partial_k \left[ (x_i-a_i)M_{ij}(x_j-a_j) \right]  \\
				     &= \frac{1}{2}\left\{ [\partial_k (x_i-a_i)]M_{ij}(x_j-a_j)+(x_i-a_i)M_{ij}[\partial_k (x_j-a_j)] \right\}  \\
				     &= \frac{1}{2}\left[ \delta_{ik}M_{ij}(x_j-a_j) + (x_i-a_i)M_{ij}\delta_{kj} \right]  \\
				     &= \frac{1}{2}\left[ M_{kj}(x_j-a_j) + (x_i-a_i)M_{ik} \right]  \\
				     &= \frac{1}{2}\left[ M_{kj}(x_j-a_j) + (x_i-a_i)M_{ki} \right]  \\
				     &= M_{kj}(x_j-a_j)
	\end{align*}
	Combining this for all components $k$, we find that
	\[
		\grad{f(\vec x)} = \mathbf{M}^2 (\vec x - \vec a)
	.\] 
	To determine the optimal step size, we again consider the vector $\vec\Delta = \vec x - \vec a$. The iteration then turns to
	\begin{align*}
		\vec{\Delta}_{k+1} &= \vec x_k - \alpha \grad{f(\vec x_k)} -\vec a\\
				   &= \vec\Delta_k - \alpha \mathbf{M}^2 \vec \Delta_k \\
				   &= (\mathbbm{1} - \alpha \mathbf{M}^2)\vec\Delta_k 
	\end{align*}
	reminding us of the previous expression in one dimension. However, now, $\mathbf{M}$ is a matrix, and thus it is not possible to tune $\alpha$ such that $\mathbbm{1}-\alpha\mathbf{M}^2=0$ unless $\mathbf{M}$ is a multiple of the identity.

	We will decompose this into a set of 1 dimensional problems by choosing the eigenbasis of $\mathbbm{1}-\alpha \mathbf{M}^2$. We will denote the eigenvalues of $\mathbbm{M}^2$ with maximum and minimum value (again, note that this is possible because $\mathbbm{M}^2$ is positive) as $\lambda_\text{max}$ and $\lambda_\text{min}$.

	For convergence, we require that $|1-\alpha \lambda|<1$ for all eigenvalues $\lambda$. This implies that $0<\lambda <2 / \alpha$. However, we only need to check the maximal eigenvalue. Thus, it converges if and only if
	 \[
		 0 < \lambda_\text{max} < \frac{2}{\alpha}
	.\] 
	The optimal $\alpha$ would minimise
	\[
	\max_i |1 - \alpha \lambda_i|
	,\]
	as the maximum value dominates the convergence. This leads to
	\[
		\alpha_\text{opt} = \frac{1}{\lambda_\text{min}+\lambda_\text{max}}
	.\] 
	Again, it is important to note that $\lambda$ are the eigenvalues of $\mathbf{M}^2$. In terms of the eigenvalues $\beta$ of $\mathbf{M}$, this would be
	\[
		\alpha_\text{opt} = \frac{1}{\beta_\text{min}^2+\beta_\text{max}^2}
	\] 
	where the subscripts max and min now denote the eigenvalues with the maximal or minimal \emph{absolute value}.\qedhere
	\end{parts}
\end{proof}
